% 5.2 — Cas d'étude : W-SEC
\section{Cas d'étude : la tâche Windows qui exécutait du code modifiable}
\label{sec:wsec}

\subsection{Le contexte}
Sous Windows, le filet est un montage WinFsp qui meurt avec la session. Une tâche planifiée le
relance à chaque ouverture de session. Le montage doit écrire dans un répertoire verrouillé au
profit des Administrateurs : la tâche s'exécute donc au niveau « le plus élevé disponible »,
une élévation accordée une fois, à l'installation, par une invite UAC que l'utilisateur a
acceptée (\fichier{schedule.py}).

\subsection{La faille}
Une tâche élevée est une cible de choix : tout ce qu'elle exécute s'exécute avec élévation.
Or, dans sa forme d'origine, la tâche lançait \cmd{\%LOCALAPPDATA\%\textbackslash
demo\_cli\textbackslash autostart-<projet>.cmd}, qui appelait \cmd{demo\_cli} trouvé par le
\cmd{PATH}. Chaque maillon était modifiable par l'utilisateur --- donc par l'agent, qui
travaille avec ses droits :

\begin{table}[htbp]
\centering
\caption{W-SEC : chaque maillon de la chaîne d'exécution, avant et après}
\label{tab:wsec}
\small
\begin{tabularx}{\linewidth}{@{}L{2.6cm}YY@{}}
\toprule
\textbf{Maillon} & \textbf{Avant (V1--V4)} & \textbf{Après} \\
\midrule
Script lancé & \cmd{.cmd} dans \cmd{\%LOCALAPPDATA\%}, modifiable (V1) &
  aucun script : la tâche lance directement l'interpréteur \\
Code exécuté & installation éditable, environnement virtuel utilisateur (V2) &
  copie figée du runtime dans \cmd{\%ProgramData\%\textbackslash demo\_cli}, verrouillé \\
Résolution & \cmd{demo\_cli} trouvé par le \cmd{PATH} de l'utilisateur (V3) &
  interpréteur \emph{de base}, chemin absolu, option \cmd{-I} (pas de \cmd{PYTHONPATH}, pas de
  site utilisateur, pas le répertoire du script) \\
Journal & écrit par la tâche élevée dans un dossier utilisateur (V4) &
  \cmd{logs\textbackslash} sous \cmd{\%ProgramData\%}, écrit par la tâche, lisible par
  l'utilisateur \\
La tâche elle-même & enregistrée par \cmd{schtasks} & enregistrée en XML avec un descripteur de
  sécurité \\
\bottomrule
\end{tabularx}
\end{table}

Les conséquences étaient concrètes : du code exécuté avec élévation à la prochaine ouverture de
session, ou immédiatement par \cmd{schtasks /run}, que l'utilisateur peut lancer sans invite ;
ou, plus discrètement, un \cmd{exit /b 0} ajouté au script pour que le garde ne revienne jamais.

\subsection{La correction}
La règle retenue (arbitrages R1--R4) tient en une phrase : \emph{rien de ce que l'utilisateur
peut écrire ne s'exécute avec élévation}.
\begin{itemize}
  \item Les fichiers de la tâche vivent dans \cmd{\%ProgramData\%\textbackslash demo\_cli},
    appartenant aux Administrateurs, verrouillé (SYSTEM et Administrateurs : contrôle total ;
    Utilisateurs : lecture et exécution). Ils sont écrits pendant l'étape déjà élevée de
    l'installation. Un dossier qui existe déjà comme lien, ou que l'utilisateur a créé avant,
    est refusé ou repris avant toute écriture.
  \item L'action de la tâche est l'interpréteur de base, avec \cmd{-I}, sur un \cmd{run.py}
    verrouillé, par chemins absolus. Un interpréteur situé dans le profil de l'utilisateur est
    refusé : il n'y aurait rien de sûr à exécuter.
  \item La tâche porte un descripteur de sécurité : les utilisateurs peuvent la lire et la
    lancer, jamais la modifier.
\end{itemize}

\begin{code}{Le descripteur de sécurité de la tâche (\fichier{schedule.py})}
# SYSTEM and Administrators full access;
# Authenticated Users read and execute (run it), never write (change it).
TASK_SDDL = "D:P(A;;FA;;;SY)(A;;FA;;;BA)(A;;GRGX;;;AU)"
\end{code}

Lecture du SDDL : \cmd{D:P} est une liste d'accès protégée (sans héritage) ; \cmd{FA} accorde le
contrôle total à SYSTEM (\cmd{SY}) et aux Administrateurs (\cmd{BA}) ; \cmd{GRGX} accorde lecture
et exécution aux utilisateurs authentifiés (\cmd{AU}).

Enfin, les répertoires de la tâche ont rejoint le plan de contrôle sous Windows, et
\cmd{demo\_cli doctor} signale une ancienne tâche, un répertoire inscriptible, un interpréteur
dans le profil de l'utilisateur ou une copie du runtime périmée.

\subsection{La vérification}
Les tests (\fichier{test\_wsec.py}) vérifient chaque décision du code sur les deux systèmes ;
ce qui exige une session Windows élevée (\cmd{schtasks /create}, \cmd{icacls} sur
\cmd{ProgramData}) a été vérifié sur Windows, en utilisateur normal : écrire dans le runtime
$\rightarrow$ refusé ; modifier la tâche $\rightarrow$ « Access is denied » ; lancer la tâche
$\rightarrow$ fonctionne.

\begin{figure}[htbp]
\centering
\capture{wsec}{La tâche durcie : refus d'écriture et de modification pour un utilisateur normal}%
{Sous Windows, en utilisateur normal : \texttt{icacls C:\textbackslash ProgramData\textbackslash demo\_cli} ; une tentative \texttt{echo x > C:\textbackslash ProgramData\textbackslash demo\_cli\textbackslash runtime\textbackslash run.py} refusée ; \texttt{schtasks /change /tn "demo\_cli guard ..." /disable} $\rightarrow$ « Access is denied ».}
\end{figure}

\begin{pourquoi}[Ce que ce cas enseigne]
La faille n'était dans aucune règle de décision : elle était dans la \emph{plomberie} qui
maintient l'outil en vie. Un garde-fou élève forcément quelque chose quelque part ; c'est là
que la revue doit regarder en premier, avec une seule question : qui peut écrire ce qui
s'exécute ?
\end{pourquoi}
